% OriginWrite starter using the generic OUP authoring class offered by G3.
% This is not the separate GSA gsag3jnl class. Check journal guidelines.
\documentclass[unnumsec,webpdf,contemporary,large,namedate]{oup-authoring-template}
\usepackage{originwrite-oup-xetex}
\newcommand{\societylogo}{}
\begin{document}
\journaltitle{G3: Genes, Genomes, Genetics}
\DOI{Manuscript draft}
\copyrightyear{2026}
\pubyear{2026}
\access{Draft for review}
\appnotes{Investigation}
\firstpage{1}
\title[Short title]{Your manuscript title}
\author[1]{Your name}
\address[1]{\orgdiv{Department}, \orgname{Institution}, \orgaddress{\country{Country}}}
\abstract{Introduce the research question, describe the approach, and summarise the main findings and their relevance.}
\keywords{genetics, genomics, methods}
\maketitle
\section{Introduction}
Introduce the biological question and its broader context.
\section{Materials and methods}
Describe the study design, materials, and analysis.
\section{Results}
Present the findings.
\section{Discussion}
Discuss implications and limitations.
\section{Data availability}
State where the supporting data and code can be accessed.
\section{Acknowledgements}
Acknowledge contributions and funding. Declare conflicts of interest.
% Link a personal library and sync references.bib before citing its papers.
% Uncomment these lines after inserting the first citation.
% \bibliographystyle{oup-abbrvnat}
% \bibliography{references}
\end{document}
